% Generated draft: requires mathematical and source review.
\documentclass[11pt]{article}
\usepackage[margin=25mm]{geometry}
\usepackage{fontspec}
\setmainfont{Latin Modern Roman}
\usepackage{amsmath,amssymb,mathtools,unicode-math}
\setmathfont{Latin Modern Math}
\usepackage{enumitem,needspace,xurl,hyperref}
\setlength{\parindent}{0pt}
\setlength{\parskip}{5pt}
\setlength{\emergencystretch}{3em}
\newcommand{\N}{\mathbb{N}}
\newcommand{\Z}{\mathbb{Z}}
\newcommand{\Q}{\mathbb{Q}}
\newcommand{\R}{\mathbb{R}}
\newcommand{\C}{\mathbb{C}}
\newcommand{\F}{\mathbb{F}}
\newcommand{\A}{\mathbb{A}}
\newcommand{\PP}{\mathbb{P}}
\newcommand{\E}{\mathbb{E}}
\newcommand{\eps}{\varepsilon}
\newcommand{\dd}{\,\mathrm{d}}
\begin{document}
\section*{UnsolvedMath ID: 30002429}
\textbf{Source problem code:} catalogue entry
\subsection{Definitions and mathematical statement}
The supplied record concerns a Gaussian limit theorem for normalized cusp-form and divisor-function values in arithmetic progressions, attributed in the supplied literature assessment to Fouvry--Ganguly--Kowalski--Michel. Let \(\Phi\) denote the auxiliary function in that theorem. The record does not specify its domain or its relationship to the progression parameters.

The proposed replacement growth condition is
\[
\exists\,\varepsilon>0\ \exists\,C_{\varepsilon}>0\ \exists\,X_{\varepsilon}>0\quad
|\Phi(x)|\leq C_{\varepsilon}x^{1-\varepsilon}
\quad\text{for every admissible }x\geq X_{\varepsilon}.
\]
Here admissible arguments and the limiting regime must be those of the referenced theorem; they are not supplied. The exponent condition uses some fixed \(\varepsilon>0\), rather than every positive \(\varepsilon\). No additional restriction such as \(\varepsilon<1\) is stated in the record.

\textbf{Requested extension, conditional on the missing theorem specification.} Does the Gaussian limit theorem for the normalized cusp-form and divisor-function values in arithmetic progressions retain its conclusion when its original growth restriction on \(\Phi\) is replaced by the displayed condition, with its other hypotheses and normalization retained?

The supplied assessment describes the established growth restriction as \(\Phi(x)=O_{\varepsilon}(x^{\varepsilon})\). It does not reproduce the theorem or explicitly state the quantification of \(\varepsilon\) in that restriction. In particular, the record does not define the random quantities, averaging measure, centering, scaling, admissible cusp forms, progression moduli, residue classes, or limiting Gaussian distribution. Consequently, this is a relative formulation of the requested extension; a precise standalone convergence statement cannot be recovered from the supplied text alone.
\subsection{Short English statement}
The problem asks whether the cited Gaussian limit theorem for cusp-form and divisor-function values in arithmetic progressions continues to hold under \(\Phi(x)=O_{\varepsilon}(x^{1-\varepsilon})\) for some fixed \(\varepsilon>0\). The original theorem's definitions and hypotheses are missing, so the intended extension requires source review before a standalone mathematical statement can be completed.
\subsection{Sources}
\begin{itemize}
\item Supplied local UnsolvedMath registry record, permanent ID \(30002429\), title ``Gaussian Limits for Arithmetic Progressions Under Weaker Growth,'' source problem code \texttt{catalogue entry}, in the Oberwolfach Reports Open Problems collection. The record gives proposed year 2013 and supplies no individual proposer.
\item Supplied report citation: \emph{Analytic Number Theory (2014)}, Oberwolfach Reports, pp.~2963--3037, DOI \texttt{10.4171/owr/2013/51}. \url{https://doi.org/10.4171/owr/2013/51}. This is also the supplied literature source labelled ``original OWR report.'' The record describes Oberwolfach Reports as documenting mathematical workshops and their open-problem discussions.
\item Supplied literature source labelled ``Fouvry--Ganguly--Kowalski--Michel'': \url{https://doi.org/10.4171/CMH/342}. The supplied assessment attributes the Gaussian theorem under \(\Phi(x)=O_{\varepsilon}(x^{\varepsilon})\) to these authors; no theorem number or full theorem statement is supplied.
\item Literature checked at: \texttt{2026-08-21}. This exact date records supplied dated curation, not independently verified research. That curation labels the problem open and describes the proposed \(O_{\varepsilon}(x^{1-\varepsilon})\) extension as conjectural in the checked literature. The background explicitly calls this a dated literature triage derived from the supplied OWR curation and requests independent verification of time-sensitive status and cited claims.
\item The registry labels the statement \texttt{corrected\_verified} and says that its mathematical content was unchanged while notation, typography, and wording were normalized into a self-contained publication form. These are supplied editorial assertions; the theorem needed for a self-contained formulation is absent from this record.
\end{itemize}
\end{document}
